%===============================================================================
% Eiffel Programming Tips
%===============================================================================
% $Id: project-description.tex 147 2005-01-24 13:41:18Z thom $
%===============================================================================


%===============================================================================
% Configuration
%===============================================================================


%-------------------------------------------------------------------------------
% \documentclass and \usepackage directives
%-------------------------------------------------------------------------------
\documentclass[a4paper,fleqn,titlepage]{article} 
%\usepackage{ngerman}
\usepackage[latin1]{inputenc}
\usepackage[T1]{fontenc}
\usepackage[small,hang,bf]{caption2}
\usepackage{fancyhdr}
\usepackage[nice]{nicefrac}
\usepackage{color,listings}
\usepackage{alltt}


% Compilation with latex or pdflatex?
\newif\ifpdf 
\ifx\pdfoutput\undefined 
  \pdffalse
\else
  \pdfoutput=1 
  \pdftrue 
\fi 

% Compilation with pdflatex
\ifpdf
 
  \usepackage[pdftex]{graphicx}

  \usepackage[
    pdftex,
    a4paper,
    bookmarks,
    pdfstartview=FitH,    % starts with page width
    bookmarksopen,        % opens index
    bookmarksnumbered,    % index with numbering
    colorlinks,           % links with color, otherwise with border
    linkcolor=blue,       % Standard red
    citecolor=blue,       % Standard green
    urlcolor=magenta,     % Standard cyan
    filecolor=blue
  ]{hyperref} 

  \pdfinfo{
    /Title      (Eiffel Programming Tips)
    /Author     (Thomas Weibel, Martin Luder, Michael K�ser)
    /Subject    (Eiffel programming)
    /Keywords   (Programming, EiffelRSS)
  }

  % Use default Acrobat reader fonts
  \usepackage{mathpazo}

  % Use CM fonts (increases document size)
  % \usepackage{ae}

% Compilation with latex
\else 

  \usepackage{graphicx} 

\fi


%-------------------------------------------------------------------------------
% Configure \maketitle
%-------------------------------------------------------------------------------
\title{EiffelRSS \\ Eiffel Programming Tips}
\author{
  Michael K\"aser <kaeserm@student.ethz.ch>
  \and 
  Martin Luder <luderm@student.ethz.ch>
  \and 
  Thomas Weibel <weibelt@student.ethz.ch>
}
\date{\today}


%-------------------------------------------------------------------------------
% Configure fancyhdr
%-------------------------------------------------------------------------------
\pagestyle{fancy}

\renewcommand{\headrulewidth}{0.1 pt}
\renewcommand{\footrulewidth}{0.1 pt}

\fancypagestyle{plain}{
  \lhead{\nouppercase{\leftmark}}
  \chead{}
  \rhead{\thepage}
  \lfoot{EiffelRSS}
  \cfoot{}
  \rfoot{Project description}
}

\lhead{\nouppercase{\leftmark}}
\chead{}
\rhead{\thepage}

\lfoot{EiffelRSS}
\cfoot{}
\rfoot{Eiffel Programming Tips}


%-------------------------------------------------------------------------------
% Configure listings
%-------------------------------------------------------------------------------
\lstset{showstringspaces=false,
  breaklines=true,
  breakindent=0pt,
  prebreak=\mbox{\tiny$\searrow$},
  postbreak=\mbox{{\color{blue}\tiny$\rightarrow$}},
  frame=trBL,
  framerule=0.75pt,
  framesep=4pt,
  rulesep=0.75pt
}


%-------------------------------------------------------------------------------
% Common configuration
%-------------------------------------------------------------------------------
\setlength{\parindent}{0em}
\setlength{\parskip}{1.5ex plus0.5ex minus0.5ex}
\sloppy
\setlength{\mathindent}{0em}


%-------------------------------------------------------------------------------
% Commandos
%-------------------------------------------------------------------------------
\newcommand{\hr}{\rule{\textwidth}{1pt}}


%===============================================================================
% Document
%===============================================================================
\begin{document}

\begin{titlepage}
  \newlength{\centeroffset}
  \setlength{\centeroffset}{-0.5\oddsidemargin}
  \addtolength{\centeroffset}{0.5\evensidemargin}

  \thispagestyle{empty}

  \noindent\includegraphics[width=\textwidth]{../figures/big_ETH}\\[-3mm]
  \hr

  \vspace*{\stretch{1}}

  \makebox[0pt][l]{
    \begin{minipage}{\textwidth}
      \flushright{
        \Huge\bfseries EiffelRSS
      }

      \noindent\rule{\textwidth}{3pt}\\[2.5ex]

      \hfill\emph{
        \Large Eiffel Programming Tips
      }
    \end{minipage}
  }

  \vspace{\stretch{1}}

  \makebox[0pt][l]{
    \begin{minipage}{\textwidth}
      \flushright{
        \bfseries 
        Michael K\"aser <kaeserm@student.ethz.ch>\\[0.3ex]
        Martin Luder <luderm@student.ethz.ch>\\[0.3ex]
        Thomas Weibel <weibelt@student.ethz.ch>\\[0.3ex]
      }
    \end{minipage}
  }

  \vspace{\stretch{1}}

  \noindent\hr\\[1mm]
  \includegraphics[width=\textwidth]{../figures/big_inf}
\end{titlepage}

% Use roman page numbering
\pagenumbering{roman}

\begin{abstract}
  This document contains various handy tips for programming with ISE
  Eiffel.
\end{abstract}

\clearpage
\tableofcontents
\newpage

% Set page counter to zero
\setcounter{page}{0} 

% Use arabic page numbering
\pagenumbering{arabic}


\section{Downloading files}
\label{sec:downloading}

With EiffelNet:

\begin{lstlisting}[language=Eiffel]
local
  http_source: HTTP_PROTOCOL
  file_destination: FILE_PROTOCOL
  http_address: HTTP_URL
  file_address: FILE_URL
  transaction: SINGLE_TRANSACTION
do
    -- create the connection to the source
  create http_address.make ("http://eiffelrss.berlios.de/Main/AllRecentChanges?action=rss")
  create http_source.make (http_address)
  
    -- print http-url-string of source
  io.put_string ("http location: ")
  io.put_string (http_source.location)
  io.put_new_line
  
    -- create the connection to the local file
  create file_address.make ("eiffelrss.rss")
  create file_destination.make (file_address)
  
    -- print file-url-sting of file
  io.put_string ("file location: ")
  io.put_string (file_destination.location)
  io.put_new_line
  
    -- create the transaction
  create transaction.make (http_source, file_destination)
  
    -- test if the transaction is correct
  io.put_string ("transaction is correct: " + transaction.is_correct.out)
  io.put_new_line
  
    -- download the file
  transaction.execute
  
    -- test if the transaction succeeded
  io.put_string ("transaction succeeded: " + transaction.succeeded.out)
  io.put_new_line
end
\end{lstlisting}

Multiple files can be downloaded with class
\texttt{MULTIPLE\_TRANSACTION}, which is a container of
\texttt{TRANSACTION}.



\section{Portable file access}
\label{sec:file}

With class \texttt{PLAIN\_TEXT\_FILE}:

\begin{lstlisting}[language=Eiffel]
local
  input_file, output_file: PLAIN_TEXT_FILE
    -- Files for input and output
do
    -- Open files
  create input_file.make_open_read ("./path/to/input.txt")
  create output_file.make_create_read_write ("./path/to/output.txt")
      
    -- Read from input
  input_file.read_line
  io.put_string (input.last_string + "%N")

    -- Write to output
  output_file.put_string (input.last_string)
  output_file.put_new_line

    -- Close files
  input_file.close
  output_file.close
end
\end{lstlisting}

This works for Unix and Windows. Make sure to only use slashes
(\texttt{/}) in the path, backslashes (\texttt{$\backslash$}) don't
work under Unix.


\section{Detecting the underlying operating system}
\label{sec:os}

The Gobo libraries feature the class \texttt{KL\_OPERATING\_SYSTEM} to
detect the underlying operating system:

\begin{lstlisting}[language=Eiffel]
local
  operating_system: KL_OPERATING_SYSTEM
do
  create operating_system
  io.put_string ("DotNet? " + operating_system.is_dotnet.out + "%N")
  io.put_string ("Unix? " + operating_system.is_unix.out + "%N")
  io.put_string ("Windows? " + operating_system.is_windows.out + "%N")
end
\end{lstlisting}


\section{Accessing features statically}
\label{sec:stat}

To access a feature from a class without creating it, just write:

\begin{lstlisting}[language=Eiffel]
feature {CLASS_NAME}.feature_name
\end{lstlisting}

For example:

\begin{lstlisting}[language=Eiffel]
create logfile.make_filename_threshold ("logfile.log", feature {LOGFILE}.Developer)
\end{lstlisting}


\section{Links}
\label{sec:links}

Links to external documentation, libraries etc.

\begin{itemize}
\item Reflection
  \footnote{http://fluri.computerscience.ch/index.php?page=projects\&subpage=m}
\item Eiffel Software Directory
  \footnote{http://eiffelzone.com/index.html}
\item ETest - Eiffel Unit Testing Tool
  \footnote{http://www.cs.yorku.ca/eiffel/etest/}
  \begin{itemize}
  \item ETest Tutorial
    \footnote{http://www.cs.yorku.ca/eiffel/prism/estudio55.htm}
  \end{itemize}
\item EiffelVision2 Tutorial
  \footnote{http://www.cs.yorku.ca/eiffel/vision2/}
\end{itemize}

\end{document}
